\chapter{ Inicio }

\section{ Primeras ideas }

Este proyecto nacio de querer entender como se crea un lenguaje, como se define
y se consume su sintaxis. De ahi surgen preguntas naturales como ¿por que
hay diferencia de velocidad entre lenguajes?

La respuesta directa es que depende de si es compilado o interpretado, pero
la pregunta real es otra: ¿por que optaron por ese diseño y como cambia el
funcionamiento del compilador?

La primera idea fue parsear con \note[expresiones regulares]{Patrones para identificar
y capturar texto} buscando estructuras tipo \backslash command\{.*\}\{.*(Contenido)\}.
Al inicio parece util, pero al anidar comandos el regex falla: el \} del comando
interior tambien cierra al superior. Se puede llevar un contador de \textbf{profundidad},
pero la expresion crece rapido y en entornos se pone peor pues ya se necesitan
funciones extra para extraer el contenido hasta el cierre.

\section{ Otras alternativas }

Probe con \note[Lark]{Herramienta para parseo y analisis gramatico} pero las 
anidaciones tampoco funcionaban de forma clara.
La siguiente idea fue trabajar caracter a caracter: acumular hasta encontrar
un caracter especial y reaccionar segun el contexto.

El problema es que el contexto complica todo. Un \{ en modo matematico es un
literal, no una accion. Eso genera demasiadas validaciones. Es como resolver
2+2 mediante union de conjuntos: mucho control para algo que podia ser simple.

\section{ Pausa }

Con esos intentos sin resultado claro, detuve el proyecto. No habia una version
\note[estable]{Una que compilara al menos el documento mas basico} lista.
Algunos estudiantes me contactaron para apoyar, les agradeci pero sin un
parser estable la estructura podia cambiar en meses y

\alert{ pudiera pasar que este proyecto se abandonara. }

Meses despues me someti a una cirugia, nada complicado pero que me brindo unos
dias de reposo y fue durante este tiempo que se me ocurrio continuar con este
proyecto, previamente abandonado.

